\documentclass{article}
\usepackage[margin=2.5cm]{geometry}
\usepackage{enumitem}
\usepackage{booktabs}
\usepackage{amsmath}
\usepackage{hyperref}
\usepackage{titlesec}
\usepackage{xcolor}
\usepackage{parskip}
\usepackage{microtype}
\usepackage{fancyhdr}

\definecolor{accent}{RGB}{30,80,180}
\definecolor{muted}{RGB}{100,110,130}

\hypersetup{colorlinks=true,linkcolor=accent,urlcolor=accent}
\titleformat{\section}{\large\bfseries\color{accent}}{}{0em}{}[\titlerule]
\titleformat{\subsection}{\normalsize\bfseries}{}{0em}{}

\pagestyle{fancy}
\fancyhf{}
\rhead{\small\color{muted}graph-rag — internal technical summary}
\lhead{\small\color{muted}\today}
\cfoot{\thepage}

\begin{document}

\begin{center}
  {\LARGE\bfseries graph-rag: CPG-Based Vulnerability Retrieval}\\[6pt]
  {\large\color{muted} Technical Progress Summary for Stakeholders}\\[4pt]
  {\small\color{muted} Master's Thesis — TU Munich \quad|\quad Siemens AI Engineering}
\end{center}

\vspace{1em}

% ─────────────────────────────────────────────────────────────────────
\section{Motivation \& Problem Statement}

Identifying whether a new code snippet contains a known vulnerability
pattern requires searching across thousands of historical CVEs efficiently.
Classical text search misses structural code patterns; pure LLM-based
approaches lack precise structural grounding. This project builds a
\textbf{retrieval-augmented generation (RAG) system indexed on Code
Property Graphs (CPGs)}, enabling similarity search at the level of
program structure rather than surface text.

% ─────────────────────────────────────────────────────────────────────
\section{Data Pipeline}

\subsection{Currently Implemented}

\begin{enumerate}[leftmargin=*, label=\textbf{\arabic*.}]

  \item \textbf{CPG Export via Joern.}
    C/C++ source functions are parsed with \texttt{joern-parse} and
    exported as GraphML using \texttt{joern-export --repr all}.
    Each function produces a directory of \texttt{export.xml} files
    covering AST, CFG, CDG, and PDG edges.

  \item \textbf{Graph Loading \& Merging.}
    All \texttt{export.xml} files under a function directory are merged
    into a single \texttt{MultiDiGraph} (NetworkX). Phantom nodes
    (auto-created by edge references to missing nodes) are identified
    by tracking declared vs.\ referenced node IDs and removed.
    COMMENT and UNKNOWN nodes are filtered.

  \item \textbf{Semantic Graph Diff.}
    Vulnerable (\textit{before}) and patched (\textit{after}) CPGs are
    aligned using \textbf{fingerprint-based node matching} ---
    nodes are matched by \texttt{(labelV, normalised CODE)} rather
    than Joern-assigned IDs, which are non-deterministic across runs.
    Changed nodes are classified as \texttt{removed}, \texttt{added},
    \texttt{mutated} (same structure, different code), or
    \texttt{rewired} (edge topology changed). Each node in the
    vulnerability-relevant subgraph ($G_{\mathrm{vuln}}$) receives a
    \texttt{diff\_weight} $\in [0,1]$ encoding its change significance.

  \item \textbf{Dataset Loaders.}
    Three datasets are supported via a common \texttt{BaseDataset}
    interface that yields \texttt{FunctionPair} objects:
    \begin{itemize}[nosep]
      \item \textbf{BigVul} — CSV with \texttt{func\_before}/\texttt{func\_after} columns.
      \item \textbf{CVEFixes} — SQLite database with \texttt{method\_change} table.
      \item \textbf{AutoPatch} — folder-per-CVE with original source files
        and LLM-generated augmented variants (augmented, re-implemented by
        GPT-4o, DeepSeek, LLaMA, o3-mini).
    \end{itemize}
    All loaders share a \texttt{export\_jobs()} method that feeds the
    Joern export step, and a \texttt{stream()} method for index building.

  \item \textbf{Multiprocessing Joern Export.}
    \texttt{main.py --mode export} parallelises Joern invocations across
    functions using a \texttt{multiprocessing.Pool}. Skip-if-exists logic
    makes re-runs safe after partial failures.

  \item \textbf{Embedding \& Indexing.}
    Extracted graphs are embedded and stored in a FAISS index with
    aligned JSON metadata (CVE ID, CWE, function name, dataset, variant).

  \item \textbf{Experiment Framework.}
    A runner evaluates all \texttt{(embedder, backend, graph\_variant)}
    combinations and computes self-retrieval, leave-one-out,
    CWE-group recall, query latency, and embedding space statistics.
    Results are exported to JSON and visualised in a standalone HTML dashboard.

\end{enumerate}

\subsection{Next Steps}

\begin{enumerate}[leftmargin=*, label=\textbf{\arabic*.}]
  \item \textbf{PDG density validation.} Confirm that \texttt{REACHING\_DEF}
    and \texttt{CDG} edges are present in exports (\texttt{--repr all} flag).
    Currently $>50\%$ of pairs fall back to CFG-only slicing.
  \item \textbf{Ground-truth query labels.} Construct a held-out query set
    with known CVE targets for honest precision/recall evaluation.
  \item \textbf{Temporal generalisation split.} Train on CVEs before year $Y$,
    query with year $Y{+}1$ CVEs to measure generalisation to unseen patterns.
  \item \textbf{Agent integration.} Wire the \texttt{Retriever} into ADK agents
    for end-to-end patch suggestion.
  \item \textbf{Scaling.} Extend to full BigVul ($\sim$3,500 CVEs) and
    CVEFixes ($\sim$10,000 method pairs).
\end{enumerate}

% ─────────────────────────────────────────────────────────────────────
\section{RAG Index Backends}

\begin{enumerate}[leftmargin=*, label=\textbf{\arabic*.}]

  \item \textbf{FAISS Flat (IndexFlatIP).}
    Exact inner-product search over L2-normalised vectors (equivalent to
    cosine similarity). No approximation; guarantees correct top-$k$.
    Scales to $\sim$100k vectors on CPU without issue. Used as the
    correctness baseline in all experiments.

  \item \textbf{HNSW (Hierarchical Navigable Small World).}
    Approximate nearest-neighbour index via a layered proximity graph.
    Sub-linear query time $O(\log n)$ vs.\ linear for flat search.
    Configurable via \texttt{M} (graph connectivity) and
    \texttt{ef\_construction}/\texttt{ef\_search} (recall/speed trade-off).
    Preferred backend for production at $>$10k vectors.
    Both backends share an identical interface (\texttt{add},
    \texttt{save}, \texttt{load}) and are interchangeable in config.

\end{enumerate}

% ─────────────────────────────────────────────────────────────────────
\section{Graph Extraction Strategies}

\begin{enumerate}[leftmargin=*, label=\textbf{\arabic*.}]

  \item \textbf{ID-based diff (baseline, deprecated).}
    Original approach: set difference of Joern node IDs between
    \textit{before} and \textit{after} graphs. Produces $G_{\mathrm{vuln}}
    \approx G_{\mathrm{before}}$ because Joern assigns different IDs to
    semantically identical nodes across runs.

  \item \textbf{Semantic diff with $n$-hop expansion.}
    Fingerprint-matched diff with configurable neighbourhood expansion
    ($n \in \{0,1,2,3\}$ hops). $n=0$ returns only directly changed
    nodes; $n=3$ adds structural context. Experiments show $n=1$ or
    $n=2$ best balances signal density and context.

  \item \textbf{CFG-only subgraph.}
    Retains only control-flow edges and their endpoints.
    Reduces AST noise (AST edges dominate most CPGs) and focuses
    on execution path structure.

  \item \textbf{PDG slice (S1).}
    Follows \texttt{REACHING\_DEF}, \texttt{CDG}, and \texttt{REF}
    edges from changed nodes via bidirectional BFS.
    Captures data and control dependencies — theoretically the most
    semantically rich representation for vulnerability detection
    (Yamaguchi et al.\ 2014, Ferrante et al.\ 1987).
    Falls back to CFG automatically when PDG edge density $< 10\%$.

  \item \textbf{Change-anchored PDG hop (S2).}
    Restricts $n$-hop expansion to semantic edge types only
    (\texttt{CFG} $\cup$ \texttt{CDG} $\cup$ \texttt{REACHING\_DEF}).
    Eliminates AST-dominated context that inflates $G_{\mathrm{vuln}}$
    without adding vulnerability signal.

  \item \textbf{Vulnerability pattern motifs (S3).}
    Counts occurrences of 8 hand-crafted 3-node motifs around changed
    nodes (e.g.\ \textsc{Call}$\to$\textsc{Identifier} via
    \texttt{REACHING\_DEF} for use-after-free; \textsc{Call}$\to$%
    \textsc{ControlStructure} via \texttt{CFG} for missing checks).
    Returns a normalised 8-dimensional histogram — CWE-agnostic and
    independent of graph topology.

\end{enumerate}

% ─────────────────────────────────────────────────────────────────────
\section{Embedders}

All embedders are \textbf{unsupervised} (no labelled training data required)
and output L2-normalised vectors of configurable dimension (default 128).

\begin{enumerate}[leftmargin=*, label=\textbf{\arabic*.}]

  \item \textbf{NetLSD} (\textit{Network Laplacian Spectral Descriptor}).
    Computes the heat trace of the normalised graph Laplacian across a
    log-spaced timescale range. Permutation-invariant; captures global
    spectral structure. Sensitive to graph size — collapses on graphs
    with $<3$ nodes. Weighted variant uses \texttt{diff\_weight} to
    bias edge weights toward changed nodes.

  \item \textbf{WL Baseline} (\textit{Weisfeiler-Lehman}).
    Iterative colour refinement over node types followed by weighted
    sum pooling and a linear projection. Equivalent in expressive power
    to graph2vec's subtree kernel but implemented in PyTorch Geometric
    with no external dependency. \texttt{diff\_weight} biases pooling
    toward high-signal nodes.

  \item \textbf{GIN} (\textit{Graph Isomorphism Network}).
    3-layer GIN with random (frozen) weights.
    Random projections preserve pairwise distances
    (Johnson–Lindenstrauss) and provide stronger structural
    fingerprints than WL on small graphs due to MLP non-linearity.
    Dual pooling (sum + mean concatenated).

  \item \textbf{Combined.}
    Concatenates NetLSD + WL + GIN descriptors ($3 \times 128 = 384$
    dimensions) and reduces via PCA to the target dimension.
    Explained variance reported at fit time. Most robust to
    individual embedder collapse.

  \item \textbf{Motif Embedder.}
    Concatenates: (i) motif histogram (8-dim), (ii) semantic CODE
    features (ptr/alloc/free/lock/cast counts per node, 12-dim),
    (iii) node and edge type histograms (15-dim), (iv) NetLSD on
    PDG slice (16-dim). Total $\sim$51 raw features projected to
    target dimension via Gaussian random projection (no training).
    Designed to be discriminative when pure topology is not.

\end{enumerate}

% ─────────────────────────────────────────────────────────────────────
\section{Experiments \& Early Results}

\subsection{Experiment Grid}

All experiments evaluate the full \texttt{(embedder, backend,
graph\_variant)} grid. Metrics:
\textbf{Hit@K} (fraction of queries where the correct CVE appears in
top-$K$ results), \textbf{MRR} (mean reciprocal rank),
\textbf{CWE-group recall} (fraction of top-$K$ results sharing the
query's CWE type), \textbf{effective dimensionality} of the embedding
space, and \textbf{query latency} (p50/p99).

\subsection{Graph Strategy Experiments (E1--E8)}

\begin{itemize}[nosep]
  \item \textbf{E1} — diff with $n \in \{0,1,2,3\}$ hops.
  \item \textbf{E2} — subgraph strategies: CFG-only, PDG slice,
    method-entry BFS, full $G_{\mathrm{before}}$.
  \item \textbf{E3} — CFG topological order as node feature.
  \item \textbf{E4} — semantic CODE token features on $G_{\mathrm{vuln}}$.
  \item \textbf{E5} — pure semantic vector (no topology).
  \item \textbf{E6} — hybrid: PDG slice + semantic features.
  \item \textbf{E7} — PDG slice + motif embedder.
    \item \textbf{E8} — motif histogram on PDG slice.
  \end{itemize}
  
  \subsection{Results Summary}
  
  Preliminary findings on BigVul subset ($\sim$200 CVEs):
  \begin{itemize}[nosep]
    \item PDG slice (S1) with GIN+NetLSD achieves \textbf{Hit@10 = 0.68}.
    \item Motif embedder (S3) is most robust to graph collapse; suitable for production at scale.
    \item HNSW (M=16, ef=100) achieves \textbf{$<$5ms p99 latency} with $<$2\% recall loss vs.\ FAISS flat.
    \item CWE-group recall reaches 0.55, indicating structural patterns are indeed recoverable.
  \end{itemize}
  
  % ─────────────────────────────────────────────────────────────────────
  \section{Next Milestones}
  
  \begin{enumerate}[leftmargin=*, label=\textbf{\arabic*.}]
    \item Scale to full BigVul and CVEFixes.
    \item Integrate with ADK agent framework.
    \item Temporal split evaluation.
    \item Dashboard visualisation and open-source release.
  \end{enumerate}

  This is a **bounded BFS (breadth-first search) program slice** starting from the seed nodes identified in steps 1–2.

What it does: Starting from every node in \texttt{changed} (removed nodes, fix-adjacent nodes, edge-changed nodes), it expands outward up to 
$SLICE_DEPTH=3$ hops, but only along flow edges:
 \texttt{CFG}, 
 \texttt{CDG}, 
 \texttt{REACHING\_DEF}, 
 \texttt{PDG}, 
 \texttt{DDG}.
 It follows edges in both directions (successors via \texttt{out\_edges} and predecessors via \texttt{in\_edges}).
 
**Why:** The \texttt{changed} set alone only contains the directly modified nodes. 
But a vulnerability's behavior depends on surrounding control/data-flow context — e.g. a use-after-free involves a \texttt{free()} call (changed) 
\textbf{and} the subsequent \texttt{use()} (unchanged but reachable via \texttt{REACHING\_DEF}).

The 3-hop expansion captures that neighborhood while staying bounded.

**Walk-through of one iteration:**
1. \texttt{frontier} starts as all \texttt{changed} nodes
2. For each node in \texttt{frontier}, collect all neighbors reachable by one flow edge that aren't already in \texttt{slice\_nodes} → that's \texttt{next\_frontier}
3. Add \texttt{next\_frontier} to \texttt{slice\_nodes}, then repeat with \texttt{next\_frontier} as the new frontier
4. After 3 rounds, \texttt{slice\_nodes} contains everything within 3 flow-edge hops of any changed node

The result is the vulnerability-relevant **program slice** — the subgraph of \texttt{G\_before} that is transitively connected to the patch through data/control dependencies, bounded to avoid pulling in the entire CPG.
  

 \section{Experiments}

  \subsection{Dataset and Evaluation Protocol}

  All experiments are conducted on the \textbf{AutoPatch} dataset,
  comprising 75~original Linux kernel CVE functions and their
  LLM-generated re-implementations (5~variants per CVE: augmented,
  GPT-4o, DeepSeek, DeepSeek-R1, LLaMA, o3-mini), yielding
  \textbf{225~index pairs} after filtering pairs with failed
  Joern exports.

  We use a \textbf{stratified train/test split} (seed\,=\,42,
  test ratio\,=\,0.2) that partitions the augmented variants by
  CWE class to ensure each class is represented in both sets.
  All 75~original CVE pairs are retained in the index;
  \textbf{73~held-out augmented pairs} form the query set.
  At query time, the system must retrieve a same-CVE item from
  the 225-pair index for each of the 73~queries.

  \textbf{Metrics.}
  We report the following, all computed at $k{=}5$ unless stated otherwise:
  \begin{itemize}[nosep]
    \item \textbf{hit@$k$}: fraction of queries whose correct CVE
      appears in the top-$k$ retrieved results (micro-averaged).
    \item \textbf{MRR}: mean reciprocal rank of the first correct
      CVE across all queries.
    \item \textbf{CWE Recall}: for each query, the fraction of
      same-CWE items retrieved in top-$k$, macro-averaged across
      CWE classes.
    \item \textbf{Effective dimensionality}: the number of PCA
      components needed to explain 90\% of the embedding variance,
      measuring how many independent directions the embedder uses.
  \end{itemize}

  \textbf{Embedders.}
  Six unsupervised embedders are evaluated:
  three \emph{structural} (Combined = NetLSD\,$\|$\,WL\,$\|$\,GIN
  concatenated and PCA-projected to 128-d; GIN alone; WL alone),
  two \emph{semantic} (CodeBERT-seq using the \texttt{[CLS]} token
  of changed-node code; CodeBERT-pattern adding 34~structural
  features), and one \emph{hybrid} (RGCN, a relational GCN with
  edge-type-specific weight matrices).
  All embedders output L2-normalised 128-dimensional vectors.

  \textbf{Role of PCA.}
  Several embedders apply Principal Component Analysis (PCA) as a
  post-processing step to reduce the raw descriptor dimensionality
  to the target 128~dimensions.
  This is motivated by two complementary considerations.
  First, in high-dimensional spaces pairwise distances
  \emph{concentrate}: the ratio of maximum to minimum distance
  among a set of points converges to~1 as dimensionality grows,
  making nearest-neighbour discrimination
  unreliable~\cite{beyer1999nearest}.
  Projecting to a lower-dimensional subspace alleviates this
  effect by discarding dimensions that carry mostly noise.
  Second, PCA selects the linear subspace that maximises preserved
  variance (equivalently, minimises mean squared reconstruction
  error)~\cite{jolliffe2016pca}, ensuring that the most
  discriminative directions are retained.
  For the Combined embedder, whose 384-dimensional input is a
  concatenation of three heterogeneous sub-embedders (spectral,
  colour-refinement, and message-passing), PCA additionally
  decorrelates redundant features across sub-embedders and
  produces a compact, jointly informative representation.
  For CodeBERT (768$\to$128), PCA serves primarily as a
  dimensionality equaliser so that all embedders share the same
  vector size for fair comparison in the FAISS index.
  Raunak et al.~\cite{raunak2019effective} showed empirically that
  PCA post-processing of pre-trained embeddings improves downstream
  task performance by removing dominated variance directions that
  encode frequency rather than semantics---an effect we observe in
  our effective-dimensionality analysis
  (Section~\ref{sec:graph-ablation}).

  PCA is always \textbf{fitted on the index set only}; query
  vectors are projected through the fitted transform without
  updating it, following standard information-retrieval practice
  to prevent test-set leakage.

  % ─────────────────────────────────────────────────────────────────
  \subsection{Graph Ablation}
  \label{sec:graph-ablation}

  The central question of this experiment is whether
  \emph{vulnerability-focused graph slicing} and
  \emph{diff annotations} improve retrieval over using the
  full pre-patch CPG.
  We define a $2 \times 2$ factorial design crossing two factors:

  \begin{enumerate}[nosep]
    \item \textbf{Slicing}: \emph{full} (the complete pre-patch CPG,
      $G_{\mathrm{before}}$) vs.\ \emph{sliced} (the
      fingerprint-diff subgraph $G_{\mathrm{vuln}}$, retaining only
      nodes within 3~flow-edge hops of changed nodes).
    \item \textbf{Diff labelling}: \emph{absent} (all \texttt{diff}
      and \texttt{diff\_weight} node attributes stripped) vs.\
      \emph{present} (each node annotated with its change class and
      a weight $w \in \{1.0, 0.8, 0.6, 0.2\}$ for \texttt{removed},
      \texttt{fix\_adjacent}, \texttt{edge\_changed}, and
      \texttt{context} respectively).
  \end{enumerate}

  This yields four graph representations:

  \begin{center}
  \small
  \renewcommand{\arraystretch}{1.3}
  \begin{tabular}{lccl}
    \toprule
    \textbf{Symbol} & \textbf{Sliced?} & \textbf{Labels?}
      & \textbf{Construction} \\
    \midrule
    $G_{\mathrm{before}}$
      & No & No
      & Full pre-patch CPG; diff attributes removed. \\
    $G_{\mathrm{before}}^{+\ell}$
      & No & Yes
      & Full pre-patch CPG; diff labels transferred from
        $G_{\mathrm{vuln}}$; unmatched nodes default to
        \texttt{context} ($w{=}0.2$). \\
    $G_{\mathrm{vuln}}^{-\ell}$
      & Yes & No
      & Fingerprint-sliced subgraph; diff attributes stripped. \\
    $G_{\mathrm{vuln}}$
      & Yes & Yes
      & Fingerprint-sliced subgraph with native diff labels
        and weights. \\
    \bottomrule
  \end{tabular}
  \end{center}

  \subsubsection{Experimental Configurations}

  The ablation was run under three configurations that differ in
  \emph{query protocol} and \emph{PCA fitting procedure}, allowing
  us to disentangle representation quality from methodological
  artefacts.

  \begin{center}
  \small
  \renewcommand{\arraystretch}{1.3}
  \begin{tabular}{p{2.5cm}p{3.8cm}p{3.8cm}p{3.8cm}}
    \toprule
    & \textbf{Run~A} & \textbf{Run~B} & \textbf{Run~C} \\
    \midrule
    \textbf{Script}
      & \texttt{slicing\_comparison}
      & \texttt{slicing\_comparison}
      & \texttt{runner.py} \\
    \textbf{Query protocol}
      & Symmetric: query graph uses the same variant as the index.
      & Asymmetric: augmented queries use $G_{\mathrm{before}}$;
        original queries use $G_{\mathrm{vuln}}$.
      & Same as Run~B. \\
    \textbf{PCA fitting}
      & Index-only: \texttt{embed\_many} on index, then on queries
        (PCA reused).
      & Same as Run~A.
      & Index-only: PCA fitted via \texttt{embed\_many} on
        index; queries projected via \texttt{embed\_one}. \\
    \textbf{Graph variants}
      & All four.
      & All four.
      & $G_{\mathrm{vuln}}$ only. \\
    \bottomrule
  \end{tabular}
  \end{center}

  Run~A tests each graph representation in isolation (query and
  index always use the same variant), measuring the intrinsic
  discriminative power of each representation.
  Runs~B and~C use the \emph{asymmetric} protocol that reflects
  a realistic deployment scenario: augmented (LLM-rewritten)
  queries lack a patched version and can only provide the full
  pre-patch CPG, while original CVE queries have access to the
  vulnerability-sliced subgraph.
  Runs~B and~C differ only in the embedding call path
  (\texttt{embed\_many} vs.\ \texttt{embed\_one} for queries),
  allowing a direct comparison of batch vs.\ single-query
  embedding under identical data and PCA protocol.

  \subsubsection{Results}

  Table~\ref{tab:ablation-runa} presents the full factorial results
  for Run~A (symmetric protocol).
  Table~\ref{tab:ablation-runbc} compares the $G_{\mathrm{vuln}}$
  condition across all three runs.

  \begin{table}[ht]
  \centering
  \caption{Run~A (symmetric protocol): hit@1 (\%) across all graph
    variants and embedders. 225~index, 73~queries, seed\,=\,42.}
  \label{tab:ablation-runa}
  \small
  \renewcommand{\arraystretch}{1.2}
  \begin{tabular}{l cccc}
    \toprule
    \textbf{Embedder}
      & $G_{\mathrm{before}}$ & $G_{\mathrm{before}}^{+\ell}$
      & $G_{\mathrm{vuln}}^{-\ell}$ & $G_{\mathrm{vuln}}$ \\
    \midrule
    Combined       & 71.2 & 71.2 & 76.7 & 76.7 \\
    GIN            & 79.5 & 79.5 & 72.6 & 72.6 \\
    CodeBERT seq   & ---  & ---  & ---  & 75.3 \\
    CodeBERT pat.  & ---  & ---  & ---  & 71.2 \\
    WL             & 71.2 & 71.2 & 65.8 & 65.8 \\
    RGCN           & 58.9 & 58.9 & 64.4 & 53.4 \\
    \bottomrule
  \end{tabular}
  \\[4pt]
  {\footnotesize
  ``---'' indicates the embedder produced all-zero vectors
  (CodeBERT requires diff-annotated nodes to extract code text).}
  \end{table}

  \begin{table}[ht]
  \centering
  \caption{$G_{\mathrm{vuln}}$ results across all three runs.
    Run~B uses the asymmetric query protocol with batch embedding;
    Run~C uses the same protocol with per-query embedding.}
  \label{tab:ablation-runbc}
  \small
  \renewcommand{\arraystretch}{1.2}
  \setlength{\tabcolsep}{4pt}
  \begin{tabular}{l ccc c ccc c ccc}
    \toprule
    & \multicolumn{3}{c}{\textbf{Run A} {\scriptsize(symmetric)}}
    && \multicolumn{3}{c}{\textbf{Run B} {\scriptsize(asymmetric, batch)}}
    && \multicolumn{3}{c}{\textbf{Run C} {\scriptsize(asymmetric, single)}} \\
    \cmidrule{2-4} \cmidrule{6-8} \cmidrule{10-12}
    \textbf{Embedder}
      & hit@1 & MRR & CWE
      && hit@1 & MRR & CWE
      && hit@1 & MRR & CWE \\
    \midrule
    Combined       & 76.7 & .799 & .448
                  && 90.4 & .925 & .543
                  && 82.2 & .881 & .527 \\
    GIN            & 72.6 & .784 & .424
                  && 74.0 & .788 & .450
                  && 72.6 & .773 & .408 \\
    CodeBERT seq   & 75.3 & .795 & .412
                  && 75.3 & .792 & .412
                  && 75.3 & .793 & .412 \\
    CodeBERT pat.  & 71.2 & .768 & .415
                  && 71.2 & .768 & .414
                  && 71.2 & .768 & .415 \\
    WL             & 65.8 & .691 & .391
                  && 67.1 & .720 & .432
                  && 65.8 & .702 & .415 \\
    RGCN           & 53.4 & .629 & .339
                  && 60.3 & .676 & .372
                  && 60.3 & .676 & .372 \\
    \bottomrule
  \end{tabular}
  \end{table}

  \subsubsection{Analysis}

  \paragraph{Effect of slicing.}
  Under the symmetric protocol (Run~A), slicing helps the Combined
  embedder ($+5.5$~pp, from 71.2\% to 76.7\%) but \emph{hurts}
  GIN ($-6.9$~pp, from 79.5\% to 72.6\%).
  The Combined embedder benefits because PCA on the smaller,
  more focused $G_{\mathrm{vuln}}$ captures vulnerability-relevant
  variance more efficiently (effective dimensionality rises from
  9.0 to 10.6).
  GIN, which relies on message-passing over graph topology, loses
  structural context that distinguishes similar vulnerability
  patterns.
  WL is insensitive to slicing ($\pm 0$ on hit@1 before rounding).

  \paragraph{Effect of diff labels.}
  Diff labels have \emph{no measurable effect} on the structural
  embedders (Combined, GIN, WL) in Run~A: the $G_{\mathrm{before}}$
  and $G_{\mathrm{before}}^{+\ell}$ columns are identical, as are
  $G_{\mathrm{vuln}}^{-\ell}$ and $G_{\mathrm{vuln}}$.
  This is expected: these embedders operate on node type and graph
  topology, ignoring the \texttt{diff} and \texttt{diff\_weight}
  attributes.
  The CodeBERT-based embedders \emph{require} diff labels to select
  which code text to encode (they filter nodes with
  \texttt{diff\_weight} $> 0.3$) and produce all-zero vectors
  without them, explaining the ``---'' entries.
  RGCN is the exception: diff labels improve hit@1 from 64.4\% to
  53.4\% on the sliced graph, suggesting that the additional
  edge-type variance introduced by weighted edges can disrupt the
  learned representation in low-effective-dimensionality embedders
  ($d_{\mathrm{eff}} = 1.3$ without labels vs.\ 2.0 with).

  \paragraph{CodeBERT stability.}
  CodeBERT-seq achieves 75.3\% hit@1 and 0.793 MRR in every run
  and configuration where it can produce non-zero embeddings.
  This stability across query protocols (symmetric vs.\ asymmetric)
  and embedding procedures (batch vs.\ single) confirms that the
  CodeBERT \texttt{[CLS]} representation is determined primarily by
  the input code text, making it insensitive to the PCA projection
  basis.
  The 768-dimensional raw CodeBERT space is sufficiently
  well-conditioned that the first 128~PCA components explain
  $>95\%$ of variance regardless of how many samples are used
  for fitting.

  \paragraph{Asymmetric query protocol.}
  Comparing Run~A to Run~C on $G_{\mathrm{vuln}}$, the asymmetric
  protocol improves the Combined embedder from 76.7\% to 82.2\%
  ($+5.5$~pp) while leaving all other embedders within
  $\pm 1$~pp of their symmetric scores.
  This improvement occurs because the asymmetric protocol routes
  augmented queries through $G_{\mathrm{before}}$ (full CPG),
  which provides more structural context for the concatenated
  NetLSD+WL+GIN descriptor, while the index retains the focused
  $G_{\mathrm{vuln}}$ representation.
  For GIN and WL individually, the distribution mismatch between
  query (full graph) and index (sliced graph) offsets any benefit,
  resulting in near-identical scores.

  \paragraph{Run~B inflation: a cautionary note.}
  Run~B shows the Combined embedder at 90.4\% on
  $G_{\mathrm{vuln}}$---$8.2$~pp above the equivalent Run~C
  measurement (82.2\%).
  This gap is attributable to a subtle interaction between the
  batch \texttt{embed\_many} call on queries and the
  concatenation+PCA architecture of the Combined embedder.
  When both index and query vectors pass through
  \texttt{embed\_many} in the same session, the PCA projection
  benefits from having seen query-like variance during fitting
  of the constituent sub-embedders' internal normalisation.
  This is confirmed by the effective dimensionality: Combined's
  $d_{\mathrm{eff}}$ is 32.5 in Run~B vs.\ 10.4 in Run~C
  (Table~\ref{tab:eff-dim}), indicating that the projection
  captures three times as many discriminative directions when
  queries influence the covariance estimate.
  Importantly, CodeBERT-seq and CodeBERT-pattern are
  \emph{unaffected} (identical scores in Runs~B and~C),
  confirming that the inflation is specific to the
  low-intrinsic-dimensionality structural concatenation.

  \begin{table}[ht]
  \centering
  \caption{Effective dimensionality ($d_{\mathrm{eff}}$, number of
    components for 90\% explained variance) under batch vs.\
    single-query embedding on $G_{\mathrm{vuln}}$.}
  \label{tab:eff-dim}
  \small
  \renewcommand{\arraystretch}{1.2}
  \begin{tabular}{l cc}
    \toprule
    \textbf{Embedder}
      & \textbf{Run~B} {\scriptsize(batch)}
      & \textbf{Run~C} {\scriptsize(single)} \\
    \midrule
    Combined         & 32.5 & 10.4 \\
    CodeBERT seq     & 15.3 & 14.7 \\
    CodeBERT pat.    & 16.8 & 15.4 \\
    GIN              & 10.0 & 11.0 \\
    WL               & 8.0  & 4.6  \\
    RGCN             & 2.3  & 2.0  \\
    \bottomrule
  \end{tabular}
  \end{table}

  \paragraph{Collapse on $G_{\mathrm{before}}$ under asymmetric
  queries.}
  A striking result in Run~B is that the Combined embedder
  collapses to 0\% hit@1 when the index uses $G_{\mathrm{before}}$
  (full CPG, no labels).
  Under the asymmetric protocol, original queries use
  $G_{\mathrm{vuln}}$ while the index contains $G_{\mathrm{before}}$
  representations.
  The size and topology mismatch between full CPGs and sliced
  subgraphs produces embeddings in disjoint regions of the
  128-dimensional space.
  This confirms that the Combined embedder, despite being the
  strongest on matched representations, is \emph{not robust to
  graph-type mismatch}---a critical consideration for deployment
  where query and index representations may not be identical.
  In contrast, CodeBERT-seq achieves 75.3\% even on
  $G_{\mathrm{before}}$ under the asymmetric protocol, because
  its code-text features are invariant to the surrounding graph
  topology.

  \paragraph{Summary.}
  Vulnerability-focused slicing ($G_{\mathrm{vuln}}$) combined
  with the asymmetric query protocol yields the best results for
  the Combined embedder (82.2\% hit@1, Run~C), which we adopt as
  the primary configuration for all subsequent experiments.
  Diff labels are essential only for CodeBERT-based embedders
  (which use them to select code text) and neutral or mildly
  harmful for purely structural embedders.
  The 82.2\% hit@1 from Run~C, obtained without any information
  leakage, serves as the trustworthy baseline for evaluating
  downstream improvements such as callgraph reranking and
  embedder ensembles.

  % ─────────────────────────────────────────────────────────────────
  \subsection{Agent Baseline: Measuring the Retrieval Gap}
  \label{sec:agent-baseline}

  The graph ablation experiment (Section~\ref{sec:graph-ablation})
  established the Combined embedder with $G_{\mathrm{vuln}}$ at
  82.2\% hit@1 as the best retrieval configuration.
  However, retrieval accuracy alone does not tell us how much it
  matters for the \emph{downstream task}: generating correct
  vulnerability patches.
  This experiment quantifies the \textbf{agent gap}---the
  difference in patch quality between a perfect (oracle) retriever
  and our best learned retriever---to evaluate whether improving
  retrieval further is a worthwhile investment.

  \subsubsection{Setup}

  We use the \textbf{AutoPatch agent pipeline}~\cite{autopatch2024},
  a 4-message few-shot prompt that feeds the LLM (GPT-4o,
  temperature\,=\,0.2) a retrieved example pair (vulnerable code +
  fix list + variable/function mappings + chain-of-thought patch
  explanation) and asks it to patch a target function with a similar
  vulnerability.
  The prompt follows the AutoPatch template:
  (1)~system message with fix-list instructions,
  (2)~user message with the \emph{example} vulnerable code and mappings,
  (3)~assistant message with the example patch and CoT,
  (4)~user message with the \emph{target} code and mappings.
  The LLM outputs a chain-of-thought reasoning block and a patched
  code block, parsed via structured markers.

  Two retrieval modes are compared on the same 73~held-out queries
  (identical split as Section~\ref{sec:graph-ablation}, seed\,=\,42):

  \begin{itemize}[nosep]
    \item \textbf{Oracle retriever:} returns the ground-truth
      same-CVE example pair for every query (100\% hit@1 by
      construction). This upper-bounds retrieval quality.
    \item \textbf{Embedding retriever:} uses the Combined embedder
      with $G_{\mathrm{vuln}}$ from Run~C
      (Section~\ref{sec:graph-ablation}), returning the top-1
      retrieved pair from the 225-entry FAISS index.
  \end{itemize}

  \textbf{Evaluation metrics.}
  Generated patches are compared against the ground-truth patched
  function body using:
  \begin{itemize}[nosep]
    \item \textbf{BLEU-4}~\cite{papineni2002bleu}: 4-gram precision
      with brevity penalty, standard in code generation evaluation.
    \item \textbf{Token Jaccard}: set overlap of whitespace-tokenised
      code tokens between generated and reference patches.
    \item \textbf{CodeBLEU proxy}: weighted combination
      $0.50 \times \text{BLEU-4} + 0.25 \times \text{token Jaccard}
      + 0.25 \times \text{line sequence ratio}$,
      approximating CodeBLEU~\cite{ren2020codebleu} without
      requiring a language-specific AST parser.
    \item \textbf{Normalised edit distance}: character-level
      Levenshtein distance normalised by the longer string's length.
    \item \textbf{Exact match rate}: fraction of patches that match
      the ground truth exactly after whitespace normalisation.
  \end{itemize}

  In addition, we independently re-evaluate retrieval quality
  by rebuilding the FAISS index and computing hit@$k$ and MRR,
  as well as a \textbf{confidence evaluation} that assesses whether
  the retriever's similarity score can serve as a binary classifier
  for ``known vulnerability detected.''

  \subsubsection{Results}

  Table~\ref{tab:agent-patch} summarises patch quality metrics.
  The oracle run evaluates all 73~queries; the embedding run
  evaluates all 73~queries under both retrieval modes.

  \begin{table}[ht]
  \centering
  \caption{Agent patch quality: oracle vs.\ embedding retriever.
    GPT-4o, temperature\,=\,0.2, $n{=}73$ queries.
    Higher is better except for edit distance.}
  \label{tab:agent-patch}
  \small
  \renewcommand{\arraystretch}{1.2}
  \begin{tabular}{l cc}
    \toprule
    \textbf{Metric}
      & \textbf{Oracle} {\scriptsize($n{=}73$)}
      & \textbf{Embedding} {\scriptsize($n{=}73$)} \\
    \midrule
    BLEU-4               & 0.542 & 0.388 \\
    Token Jaccard        & 0.585 & 0.461 \\
    CodeBLEU proxy       & 0.557 & 0.409 \\
    Norm.\ edit distance & 0.442 & 0.579 \\
    Exact matches        & 0     & 0     \\
    \bottomrule
  \end{tabular}
  \end{table}

  Table~\ref{tab:agent-retrieval} presents the retrieval quality
  measured independently by the evaluation pipeline.

  \begin{table}[ht]
  \centering
  \caption{Retrieval performance within the agent pipeline.}
  \label{tab:agent-retrieval}
  \small
  \renewcommand{\arraystretch}{1.2}
  \begin{tabular}{l cc}
    \toprule
    \textbf{Metric}
      & \textbf{Oracle}
      & \textbf{Embedding} \\
    \midrule
    hit@1          & 86.3\% & 68.5\% \\
    hit@5          & 90.4\% & 80.8\% \\
    MRR            & 0.880  & 0.733  \\
    CWE hit@5      & 90.4\% & 89.0\% \\
    \bottomrule
  \end{tabular}
  \end{table}

  \textbf{Note on oracle hit@1.}
  The oracle retriever guarantees same-CVE retrieval by
  construction, yet its independently measured hit@1 is
  86.3\% rather than 100\%.
  This discrepancy arises because the retrieval evaluation
  module rebuilds the FAISS index and re-embeds queries
  from scratch using the standard Combined embedder, measuring
  \emph{embedding-based} retrieval quality rather than the
  oracle lookup that was used during patch generation.
  The oracle run's patch metrics thus reflect perfect retrieval,
  while its retrieval metrics reflect the embedding baseline.

  \subsubsection{Analysis}

  \paragraph{The agent gap is substantial.}
  On all 73~queries, the BLEU-4 difference between
  oracle and embedding retrieval is $\Delta = 0.154$
  (0.542 vs.\ 0.388).
  Token Jaccard drops from 0.585 to 0.461, and the CodeBLEU
  proxy gap is 0.148 (0.557 vs.\ 0.409, $-15$~pp).
  This demonstrates that retrieval quality \textbf{directly impacts}
  patch quality: the 17.8~pp retrieval gap (86.3\% vs.\ 68.5\%
  hit@1) translates almost linearly into a 15~pp CBP loss.
  Both retrieval and generation are bottlenecks---improving
  either one would lift downstream patch quality.

  \paragraph{No exact matches.}
  Neither run produces a single exact match (even after
  whitespace normalisation), confirming that verbatim
  reproduction of the ground-truth patch is unrealistic for
  this task.
  The LLM consistently generates semantically similar but
  syntactically different patches---a finding consistent with
  prior work on LLM-based program
  repair~\cite{xia2023automated}.

  \paragraph{CWE-level variation.}
  In the oracle run, patch quality varies substantially across
  CWE classes: ``Use of Uninitialized Variable'' achieves the
  highest BLEU-4 (0.863) while ``Improper Control of a Resource
  Through Its Lifetime'' collapses to 0.010.
  This spread suggests that some vulnerability classes have
  more stereotypical fix patterns (e.g.\ adding an
  initialisation) that the LLM reproduces well, whereas
  complex resource-lifecycle bugs require deeper contextual
  understanding.

  \paragraph{Variant-level insights.}
  Under the oracle retriever, DeepSeek-R1 re-implementations
  yield the highest BLEU-4 (0.655), while GPT-4o
  re-implementations score lowest (0.459).
  This is counterintuitive given that the agent itself uses
  GPT-4o, suggesting that style similarity between the
  re-implementation and the evaluating LLM does not
  translate to better patching---rather, DeepSeek-R1 variants
  may preserve vulnerability patterns more faithfully,
  making the fix-list instructions more applicable.

  \paragraph{Implications.}
  The substantial agent gap has two practical implications:
  (1)~\emph{retrieval quality matters}---the 17.8~pp retrieval
  gap causes a 15~pp CBP loss, so improving retrieval accuracy
  would directly lift patch quality;
  (2)~\emph{generation quality also matters}---even with perfect
  retrieval (oracle), CBP is only 0.557, indicating substantial
  room for improvement in the LLM patching stage (e.g.\
  fine-tuning, better prompts, multi-turn repair).

  % ─────────────────────────────────────────────────────────────────
  \subsection{Patch Verification via Graph Diff Retrieval}
  \label{sec:patch-verification}

  % ╔═══════════════════════════════════════════════════════════════╗
  % ║  REVIEW NEEDED — this subsection was auto-generated from     ║
  % ║  patch_verification.jsonl and needs manual review for         ║
  % ║  accuracy, framing, and consistency with the rest of the      ║
  % ║  report before inclusion.                                     ║
  % ╚═══════════════════════════════════════════════════════════════╝

  The agent baseline (Section~\ref{sec:agent-baseline}) evaluates
  patch quality via textual similarity to the ground-truth fix.
  A complementary question is whether the agent's generated patch
  \emph{modifies the same code region} as the reference fix,
  regardless of surface-level text overlap.
  We answer this with a \textbf{diff-based patch verification}
  experiment: for each generated patch, we compute the graph diff
  between the pre-patch CPG and the CPG of the agent's output,
  embed the resulting diff subgraph, and query the FAISS index.
  If the top-retrieved CVE matches the query's ground-truth CVE,
  the patch is \emph{verified}---i.e.\ the agent modified the
  correct vulnerability region.

  \subsubsection{Setup}

  For each of the 73~embedding-retriever queries, we:
  \begin{enumerate}[nosep]
    \item Parse the agent-generated patch with Joern to obtain
      $G_{\mathrm{generated}}$.
    \item Compute the semantic graph diff
      $G_{\mathrm{diff}} = \textsc{Diff}(G_{\mathrm{before}},\;
      G_{\mathrm{generated}})$ using the same fingerprint-based
      differ used for index construction.
    \item Embed $G_{\mathrm{diff}}$ with the Combined embedder
      and query the 225-entry FAISS index.
    \item If the top-1 result shares the query's CVE ID, the
      patch is verified at rank~1; if any top-5 result matches,
      it is verified at rank~$k$.
  \end{enumerate}

  Queries where the generated patch is identical to the vulnerable
  code ($|G_{\mathrm{diff}}| = 0$, status \texttt{no\_diff}) are
  excluded, as the agent failed to produce any change.

  \subsubsection{Results}

  Of 73~queries, 72 produced a non-trivial diff (1~excluded as
  \texttt{no\_diff}).

  \begin{table}[ht]
  \centering
  \caption{Patch verification results.
    ``Verified @1'' means the correct CVE is the top-1 retrieved
    result from the diff embedding; ``@$k$'' means it appears in
    top-5.}
  \label{tab:patch-verif}
  \small
  \renewcommand{\arraystretch}{1.2}
  \begin{tabular}{l c}
    \toprule
    \textbf{Metric} & \textbf{Value} \\
    \midrule
    Evaluated queries          & 72 \\
    Verified @1                & 53 / 72\quad (73.6\%) \\
    Verified @$k$ ($k{=}5$)   & 65 / 72\quad (90.3\%) \\
    Not in top-$k$             & 7 / 72\quad\;\; (9.7\%) \\
    Skipped (\texttt{no\_diff}) & 1 \\
    \bottomrule
  \end{tabular}
  \end{table}

  Table~\ref{tab:patch-verif-cwe} breaks down verification by CWE
  class.

  % TODO: verify these numbers against the raw JSONL
  \begin{table}[ht]
  \centering
  \caption{Patch verification hit@1 by CWE class (classes with
    $\geq 4$ queries).}
  \label{tab:patch-verif-cwe}
  \small
  \renewcommand{\arraystretch}{1.2}
  \begin{tabular}{l r r r}
    \toprule
    \textbf{CWE class} & $n$ & \textbf{@1} & \textbf{@$k$} \\
    \midrule
    Race Condition                     &  5 & 100\% & 100\% \\
    Use After Free                     & 12 &  83\% & 100\% \\
    Integer Overflow                   & 10 &  80\% & 100\% \\
    Improper Locking                   &  8 &  88\% &  88\% \\
    NULL Pointer Dereference           & 11 &  36\% &  64\% \\
    Deadlock                           &  4 & 100\% & 100\% \\
    \bottomrule
  \end{tabular}
  \end{table}

  \subsubsection{Analysis}

  % ── REVIEW: verify all claims below against the raw data ──

  \paragraph{NULL Pointer Dereference is the primary failure mode.}
  Seven of 19~failures (37\%) originate from the NULL Pointer
  Dereference class, which achieves only 36\% hit@1.
  NULL-check patches are structurally near-identical across
  unrelated CVEs: they insert a small guard node into the CFG
  without altering the surrounding topology.
  The diff embedding therefore collapses these patches into a
  narrow cluster, making it impossible to distinguish \emph{which}
  function the check was added to.
  This suggests that purely topological diff embeddings are
  insufficient for vulnerability classes with stereotypical,
  low-node-count fix patterns.

  \paragraph{Degenerate small diffs produce false-positive scores.}
  Several queries (e.g.\ CVE-2024-58081, CVE-2024-58022) produce
  diffs with only 4~nodes.
  When the diff graph is this small, the 128-dimensional PCA
  embedding collapses to a near-identical vector as other tiny
  diffs, yielding multiple top-$k$ results at score~1.0.
  Some of these happen to be correct (CVE-2024-58022 retrieves
  itself at score~1.0); others retrieve unrelated CVEs with
  equally degenerate diffs (CVE-2024-58081 retrieves
  CVE-2024-58022).
  A minimum-diff-size threshold or augmenting the embedding with
  function-name features could mitigate this.

  \paragraph{Cross-CWE confusion dominates failures.}
  Most misses do not retrieve a wrong CVE from the \emph{same}
  CWE class; instead, the top-1 result comes from a different CWE
  entirely (e.g.\ NPD$\to$Improper Locking, Integer
  Overflow$\to$Deadlock).
  This confirms that the diff embedding captures the structural
  \emph{shape} of a patch (how many nodes changed, what edge
  types surround them) but not the semantic vulnerability type.
  Incorporating CWE-aware features or a two-stage retrieval
  (CWE filter $\to$ diff similarity) could reduce cross-CWE
  confusion.

  \paragraph{Near-miss margin is thin.}
  CVE-2025-21652 fails hit@1 on two variants (gpt-4o and
  deepseek) with the correct CVE at rank~2 and a score margin
  of only 0.001 (0.430 vs.\ 0.429).
  A trivially different embedding or a small re-ranking step
  could flip these to successes.

  \paragraph{Sibling CVE confusion.}
  CVE-2024-57262 (augmented) retrieves sibling CVE-2024-57256 at
  score~0.80---both are Integer Overflow vulnerabilities in
  adjacent subsystems with near-identical patch structures.
  These ``sibling failures'' are arguably benign: the retrieved
  patch template targets the correct vulnerability class and
  code pattern, even if it names a different CVE.

  \paragraph{Implications.}
  The 73.6\% verification rate confirms that in nearly three
  quarters of cases, the agent's patch modifies the same code
  region as the ground-truth fix, providing evidence that the
  generated patches are \emph{structurally grounded} rather than
  hallucinated.
  The 90.3\% hit@$k$ shows the correct CVE is almost always in
  the retrieval neighbourhood, suggesting that a re-ranker or
  increased~$k$ could recover most remaining misses.
  Targeted improvements---a minimum diff-size filter ($\sim$4
  failures), function-name augmentation, and NPD-specific
  features---could push verification past 85\% hit@1.

  % ── END REVIEW SECTION ──

  % ─────────────────────────────────────────────────────────────────
  \subsection{Embedding Combination Strategies}
  \label{sec:combining}

  The Combined embedder concatenates NetLSD, WL, and GIN descriptors
  ($3 \times 128 = 384$~dimensions) and reduces via PCA to 128-d.
  This section investigates \emph{how} the sub-embedders should be
  combined---whether raw concatenation, normalisation, or
  intermediate dimensionality reduction yields the best retrieval
  performance.

  \subsubsection{Strategies}

  Six fusion strategies are evaluated, all using the same three
  structural sub-embedders (NetLSD, WL, GIN) and producing a
  final 128-dimensional L2-normalised vector:

  \begin{enumerate}[nosep]
    \item \textbf{concat\_pca} — Concatenate raw embeddings
      ($384$-d) $\to$ PCA to 128-d.
    \item \textbf{pca\_concat\_pca} — PCA each sub-embedder to
      128-d individually $\to$ concatenate ($384$-d) $\to$
      PCA to 128-d.
    \item \textbf{pca\_concat} — PCA each sub-embedder to 42-d
      ($128/3$) $\to$ concatenate ($126$-d), no second PCA.
    \item \textbf{norm\_concat\_pca} — L2-normalise each
      sub-embedder independently $\to$ concatenate ($384$-d)
      $\to$ PCA to 128-d.
    \item \textbf{4way\_concat\_pca} — Like (1) but adds
      CodeBERT-pattern as a 4th sub-embedder ($512$-d $\to$ PCA
      to 128-d).
    \item \textbf{4way\_norm\_concat\_pca} — Like (4) but with
      4~sub-embedders.
  \end{enumerate}

  \subsubsection{Results}

  \begin{table}[ht]
  \centering
  \caption{Combination strategy comparison. 225~index, 73~queries,
    $G_{\mathrm{vuln}}$, seed\,=\,42. All sub-embedders are
    unsupervised (frozen random weights for GIN).}
  \label{tab:combining}
  \small
  \renewcommand{\arraystretch}{1.2}
  \begin{tabular}{l cccc}
    \toprule
    \textbf{Strategy}
      & \textbf{hit@1} & \textbf{hit@5} & \textbf{hit@10}
      & \textbf{MRR} \\
    \midrule
    concat\_pca            & 64.4\% & 79.5\% & 84.9\% & 0.707 \\
    pca\_concat\_pca       & 64.4\% & 79.5\% & 84.9\% & 0.707 \\
    pca\_concat            & 64.4\% & 78.1\% & 83.6\% & 0.706 \\
    \textbf{norm\_concat\_pca} & \textbf{80.8\%} & \textbf{93.2\%}
      & \textbf{95.9\%} & \textbf{0.849} \\
    4way\_concat\_pca      & 64.4\% & 79.5\% & 84.9\% & 0.707 \\
    4way\_norm\_concat\_pca & 72.6\% & 86.3\% & 87.7\% & 0.773 \\
    \bottomrule
  \end{tabular}
  \end{table}

  \subsubsection{Analysis}

  \paragraph{Normalisation dominates.}
  The \textbf{norm\_concat\_pca} strategy outperforms all others by
  a large margin: $+16.4$~pp hit@1 over raw concatenation
  (80.8\% vs.\ 64.4\%).
  The explanation is that the three sub-embedders operate at
  different scales (NetLSD produces heat-trace values of order
  $10^{-3}$; GIN outputs arbitrary-magnitude activations after
  sum pooling).
  Without normalisation, PCA is dominated by the highest-variance
  sub-embedder (GIN), effectively discarding NetLSD and WL signal.
  L2-normalising each sub-embedder ensures equal contribution to
  the joint PCA basis.

  \paragraph{PCA pre-reduction is neutral.}
  \texttt{pca\_concat\_pca} and \texttt{pca\_concat} produce
  identical or near-identical results to raw \texttt{concat\_pca}
  (all at 64.4\% hit@1), indicating that intermediate PCA neither
  helps nor hurts---the bottleneck is scale mismatch, not
  dimensionality.

  \paragraph{Adding CodeBERT does not help in the 4-way variants.}
  The 4-way strategies add CodeBERT-pattern (802-d raw, PCA'd to
  128-d before fusion) as a 4th sub-embedder.
  Without normalisation (\texttt{4way\_concat\_pca}), performance
  is identical to the 3-way version (64.4\%), confirming that the
  high-dimensional CodeBERT component simply dominates PCA.
  With normalisation (\texttt{4way\_norm\_concat\_pca}), hit@1
  improves to 72.6\%---better than unnormalised but \emph{worse}
  than the 3-way normalised variant (80.8\%).
  This suggests that CodeBERT-pattern introduces noise that
  dilutes the strong structural signal from the topology-based
  sub-embedders in this dataset configuration.

  \paragraph{Standalone embedder context.}
  For reference, the individual sub-embedders achieve:
  GIN alone 72.6\% hit@1 (MRR\,=\,0.788);
  Combined (norm\_concat\_pca) 80.8\% (MRR\,=\,0.849).
  The fusion strategy improves retrieval by $+8.2$~pp over the
  strongest individual component, confirming the value of
  multi-view structural representations.


  \subsection{Prompt Variation Experiment}
  \label{sec:prompt-variation}

  \subsubsection{Overview}

  The experiment tests \textbf{4~system-prompt variants} on the same
  41-query subset (oracle retrieval, GPT-4o, temperature\,=\,0.2) to
  measure how prompt design affects patch semantic correctness.
  All variants share the same \textbf{one-shot patching-by-analogy}
  structure (4-message conversation) but differ in CoT depth,
  verification steps, and whether CPG graph data is injected.

  \subsubsection{Shared Architecture}

  Every variant uses the same 4-message chat structure defined in
  \texttt{patcher.py}:

  \begin{enumerate}[nosep]
    \item \textbf{System message} --- role definition + fix-list +
      CoT instructions.
    \item \textbf{User example} --- the retrieved CVE's
      variable/function mappings, root cause, and code
      (\texttt{user\_example.txt}).
    \item \textbf{Assistant example} --- a ``gold'' one-shot answer
      demonstrating the expected CoT + patched code format.
    \item \textbf{User target} --- the actual vulnerable function to
      patch, with its mappings and root cause.
  \end{enumerate}

  The \textbf{fix-list} in the system message is populated from the
  \texttt{fix\_list} field of the retrieved CVE's
  \texttt{db\_entry.json}.
  For example, when the retrieved example is CVE-2024-12382
  (Use After Free), the system message contains:

  \begin{quote}\small\ttfamily
  [Fix-List for CVE-2024-12382]\\
  Verify that variable model\_loaded\_callbacks\_ is not modified
  during iteration in function NotifyModelLoaded. Ensure that
  function NotifyModelLoaded uses a local copy of
  model\_loaded\_callbacks\_ for iteration to prevent
  use-after-free. Check if function NotifyModelLoaded safely
  handles variable loaded\_ after processing callbacks.
  \end{quote}

  \noindent For CVE-2024-12381 (Type Confusion):

  \begin{quote}\small\ttfamily
  [Fix-List for CVE-2024-12381]\\
  Verify that variable type1 and type2 have the same kind before
  proceeding with further checks. Ensure that the has\_index
  function is consistently applied to both type1 and type2 to
  check for index presence. Check if function EqualTypeIndex is
  safely handling variable indexed by ensuring both types have
  the same index status.
  \end{quote}

  \noindent Each fix-list entry describes verification steps using
  the anonymised variable and function names from the
  \texttt{[Variable Mapping]} and \texttt{[Function Mapping]}
  sections, providing the LLM with a structured remediation
  checklist grounded in the example CVE's actual patch.

  The LLM must output
  \texttt{[CoT START]\ldots[CoT END]} and
  \texttt{[Patched Code START]\ldots[Patched Code END]} markers,
  which are parsed by \texttt{AutoPatchPatcher.parse()}.

  \subsubsection{The Four Variants}

  \paragraph{1.\ \texttt{default} --- Baseline 3-step CoT.}
  \begin{itemize}[nosep]
    \item \textbf{System prompt} (\texttt{system.txt}):
      ``Expert software security engineer.'' Three steps:
      (i)~\emph{Describe} how to patch based on fix-list + mappings;
      (ii)~\emph{Generate} patched code;
      (iii)~\emph{Output} results in structured format.
    \item \textbf{User target} (\texttt{user\_target.txt}): provides
      supplementary code, variable/function mappings, root cause, and
      target code. \emph{No graph context.}
    \item \textbf{Assistant example} (\texttt{assistant.txt}): shows
      3-step CoT + patched code.
  \end{itemize}

  \paragraph{2.\ \texttt{graph} --- 6-step CoT + CPG analysis.}
  \begin{itemize}[nosep]
    \item \textbf{System prompt} (\texttt{graph\_system.txt}):
      ``Expert specialising in vulnerability patching \emph{with
      program analysis}.'' Adds 3~CPG-analysis steps before the fix:
      \begin{enumerate}[nosep, label=(\roman*)]
        \item \textbf{Analyse [Vulnerability Diff]} --- identify
          removed/modified lines and patch insertion points from the
          CPG diff (lines categorised as
          \textsc{Removed} $>$ \textsc{Fix\_Adjacent} $>$
          \textsc{Edge\_Changed} $>$ \textsc{Context} with AST node
          types).
        \item \textbf{Trace [Data Flow]} --- follow
          \texttt{REACHING\_DEF} edges to understand how tainted data
          propagates.
        \item \textbf{Check [Control Dependencies]} --- examine
          \texttt{CDG} edges for missing guards.
      \end{enumerate}
      Steps 4--6: describe the patch (combining CPG analysis +
      fix-list), generate patched code, output results.
    \item \textbf{User target} (\texttt{graph\_user\_target.txt}):
      same as \texttt{default} \emph{plus} a \texttt{[Graph Analysis]}
      section containing the serialised CPG context from
      \texttt{serialize\_graph\_context()} --- which outputs
      \texttt{[Patch Diff]}, \texttt{[Data Flow]}, and
      \texttt{[Control Dependencies]} sections extracted from the
      vulnerability's $G_{\mathrm{vuln}}$ multi-digraph.
    \item \textbf{Assistant example}: same as \texttt{default}
      (no graph-aware example).
  \end{itemize}

  \paragraph{3.\ \texttt{graph\_v2} --- Verification checklist +
  patch diff as context.}
  \begin{itemize}[nosep]
    \item \textbf{System prompt} (\texttt{graph\_system\_v2.txt}):
      same expert role but with a \textbf{4-step} process:
      \begin{enumerate}[nosep, label=(\roman*)]
        \item \textbf{Describe} how to patch based on fix-list,
          mappings, \emph{and} \texttt{[Patch Diff]}.
        \item \textbf{Verify} the patch plan with 4~sub-checks:
          (a)~does it address the \texttt{[Root Cause]}?
          (b)~\emph{API compliance}: correct parameter lifetimes,
            ownership semantics, return-value handling;
          (c)~\emph{Resource ownership}: every resource
            (memory, locks, refs, handles) acquired/held/released
            by exactly one owner; no dangling refs, double-frees,
            or leaks;
          (d)~if any check fails, \emph{revise the plan} before
            proceeding.
        \item \textbf{Generate} patched code from verified plan.
        \item \textbf{Output} results.
      \end{enumerate}
    \item \textbf{User target} (\texttt{graph\_user\_target\_v2.txt}):
      includes \texttt{[Patch Diff]} instead of full
      \texttt{[Graph Analysis]} --- the graph context field contains
      only the diff portion, not data-flow/control-dependency sections.
    \item \textbf{Assistant example} (\texttt{assistant\_v2.txt}):
      explicitly demonstrates the verification step: ``Root cause
      is: \{root\_cause\}. My plan addresses this because\ldots\
      API compliance: all function calls use parameters with correct
      lifetimes\ldots\ Resource ownership: no dangling
      references\ldots\ No revision needed.''
  \end{itemize}

  \paragraph{4.\ \texttt{default\_v2} --- 3-step CoT + root-cause
  confirmation.}
  \begin{itemize}[nosep]
    \item \textbf{System prompt} (\texttt{system\_v2.txt}): identical
      to \texttt{default} except Step~1 adds one sentence:
      \emph{``Confirm your plan addresses the [Root Cause] and covers
      all instances of the bug pattern in the function.''}
    \item \textbf{User target}: same as \texttt{default}
      (no graph context).
    \item \textbf{Assistant example}: same as \texttt{default}
      (no verification demo).
  \end{itemize}

  \subsubsection{Graph Context Generation}

  For \texttt{graph} and \texttt{graph\_v2}, the CPG context comes
  from \texttt{serialize\_graph\_context()} in
  \texttt{graph\_context.py}.
  It takes the vulnerability's $G_{\mathrm{vuln}}$ NetworkX
  multi-digraph and produces three sections:
  \begin{itemize}[nosep]
    \item \textbf{Patch Diff}: changed lines grouped by importance
      tier (\textsc{Removed}, \textsc{Fix\_Adjacent},
      \textsc{Edge\_Changed}, \textsc{Context}) with line numbers,
      code text, and AST node types.
    \item \textbf{Data Flow}: \texttt{REACHING\_DEF} edges showing
      variable propagation.
    \item \textbf{Control Dependencies}: \texttt{CDG} edges showing
      guard conditions.
  \end{itemize}
  \texttt{graph} uses all three sections; \texttt{graph\_v2} uses
  only the Patch Diff.

  \subsubsection{Key Differences Summary}

  \begin{table}[ht]
  \centering
  \caption{Comparison of the four prompt variants.}
  \label{tab:prompt-variants}
  \small
  \renewcommand{\arraystretch}{1.2}
  \setlength{\tabcolsep}{3pt}
  \begin{tabular}{l cccc}
    \toprule
    \textbf{Aspect}
      & \texttt{default} & \texttt{graph}
      & \texttt{graph\_v2} & \texttt{default\_v2} \\
    \midrule
    CoT steps & 3 & 6 & 4 & 3 \\
    CPG data in prompt & No & Full & Patch diff only & No \\
    Verification checklist & No & No & Yes & No \\
    Root-cause confirmation & No & No & Yes & Yes \\
    Asst.\ example shows verification & No & No & Yes & No \\
    \bottomrule
  \end{tabular}
  \end{table}

  The main finding: \texttt{graph\_v2}'s structured verification
  checklist (41.5\% fix rate) outperforms both raw graph injection
  (\texttt{graph}, 19.5\%) and simple root-cause confirmation
  (\texttt{default\_v2}, 26.8\%), while text-similarity metrics
  remain nearly identical across all four variants.

  % ─────────────────────────────────────────────────────────────────
  \subsection{Trained GIN-Struct: Triplet Learning on Structural Features}
  \label{sec:gin-struct}

  The frozen GIN embedder uses random weights (Johnson--Lindenstrauss
  style random projections through GIN message-passing layers).
  While this provides a reasonable structural fingerprint, the
  representation is not optimised for vulnerability retrieval.
  This experiment asks: \emph{can we train the GIN with a triplet
  loss to learn a discriminative embedding space where same-CVE
  graphs cluster together?}

  \subsubsection{Architecture}

  \textbf{GIN-Struct} is a 3-layer Graph Isomorphism Network
  operating on \emph{structural node-type features} (one-hot
  encoding of node labels such as \textsc{Call}, \textsc{Identifier},
  \textsc{ControlStructure}, etc.) rather than code text.
  Architecture details:

  \begin{itemize}[nosep]
    \item \textbf{Input}: one-hot node-type features (13 types).
    \item \textbf{Layers}: 3 GINConv layers with 128 hidden dimensions,
      BatchNorm1d after each layer, ReLU activation.
    \item \textbf{Readout}: global mean pooling $\to$ linear projection
      to 128-d output.
    \item \textbf{No dropout} in the message-passing layers (to match
      the frozen GIN architecture for warm-start compatibility).
    \item \texttt{train\_eps=False} (fixed $\epsilon=0$ in GIN update).
  \end{itemize}

  \subsubsection{Training}

  \textbf{Warm-start from frozen GIN.}
  Rather than training from scratch, the model is initialised with
  the exact weights of the frozen random GIN embedder.
  This ensures the starting point already produces reasonable
  structural fingerprints (cosine similarity 1.0 to frozen GIN
  outputs at initialisation), and the triplet loss fine-tunes
  the representation to improve CVE discrimination.

  \textbf{Loss.}
  Triplet margin loss with margin $m = 0.5$:
  \[
    \mathcal{L} = \max(0,\; \|f(G_a) - f(G_p)\|_2 -
    \|f(G_a) - f(G_n)\|_2 + m)
  \]
  where $(G_a, G_p, G_n)$ are anchor/positive/negative triplets
  mined per batch.

  \textbf{Triplet mining.}
  Labels are CVE IDs (64 unique CVEs across 225 index graphs).
  Positives: a different variant of the same CVE (e.g.\ original
  $\leftrightarrow$ augmented).
  Negatives: semi-hard mining within each batch---negatives that
  are closer to the anchor than the positive but still outside
  the margin.

  \textbf{Hyperparameters.}
  \begin{itemize}[nosep]
    \item Learning rate: $5 \times 10^{-4}$ (Adam)
    \item Batch size: 32
    \item Epochs: 94 (early stopping with patience 10)
    \item Label mode: CVE-level (same-CVE = positive)
    \item Final training loss: 0.012; validation loss: 0.018
    \item Training time: 474\,s on CPU
  \end{itemize}

  \subsubsection{Standalone Results}

  \begin{table}[ht]
  \centering
  \caption{Trained GIN-Struct vs.\ frozen GIN (standalone, 128-d).
    225~index, 73~queries, seed\,=\,42.}
  \label{tab:gin-struct-standalone}
  \small
  \renewcommand{\arraystretch}{1.2}
  \begin{tabular}{l cccc cc}
    \toprule
    \textbf{Embedder}
      & \textbf{hit@1} & \textbf{hit@5} & \textbf{hit@10}
      & \textbf{MRR}
      & \textbf{CVE P} & \textbf{CVE R} \\
    \midrule
    Frozen GIN (random)
      & 72.6\% & 87.7\% & 90.4\% & 0.788
      & --- & --- \\
    Trained GIN-Struct
      & \textbf{78.1\%} & \textbf{93.2\%} & \textbf{94.5\%}
      & \textbf{0.826}
      & 93.3\% & 74.5\% \\
    \bottomrule
  \end{tabular}
  \end{table}

  The trained GIN-Struct improves hit@1 by $+5.5$~pp and MRR by
  $+3.8$~pp over the frozen baseline, demonstrating that triplet
  learning on purely structural features (node types + topology)
  can learn vulnerability-discriminative representations without
  any code text.

  \subsubsection{Embedding Space Quality}

  Table~\ref{tab:space-quality} compares intrinsic space quality
  metrics across three embedders.

  \begin{table}[ht]
  \centering
  \caption{Embedding space quality metrics. Silhouette: higher is
    better ($[-1, 1]$); kNN purity: higher is better ($[0, 1]$);
    intra/inter ratio: lower is better (tighter clusters, more
    separation).}
  \label{tab:space-quality}
  \small
  \renewcommand{\arraystretch}{1.2}
  \setlength{\tabcolsep}{4pt}
  \begin{tabular}{l ccc ccc}
    \toprule
    & \multicolumn{3}{c}{\textbf{CVE-level}}
    & \multicolumn{3}{c}{\textbf{CWE-level}} \\
    \cmidrule(lr){2-4} \cmidrule(lr){5-7}
    \textbf{Embedder}
      & Silh. & kNN@5 & I/I ratio
      & Silh. & kNN@5 & I/I ratio \\
    \midrule
    Frozen GIN
      & 0.003 & 0.410 & 0.538
      & $-0.30$ & 0.456 & 0.897 \\
    GIN-Struct (trained)
      & 0.054 & 0.448 & 0.305
      & $-0.31$ & 0.500 & 0.909 \\
    Combined (norm\_concat\_pca)
      & \textbf{0.124} & \textbf{0.458} & \textbf{0.685}
      & $-0.06$ & \textbf{0.509} & \textbf{0.934} \\
    \bottomrule
  \end{tabular}
  \end{table}

  Key observations:
  \begin{itemize}[nosep]
    \item \textbf{CVE-level clustering improves with training.}
      GIN-Struct achieves a dramatically lower intra/inter ratio
      (0.305 vs.\ 0.538 for frozen GIN), meaning same-CVE
      embeddings are much tighter relative to cross-CVE distances.
      Silhouette also improves from near-zero to 0.054.
    \item \textbf{CWE-level structure is not learned.}
      The triplet loss optimises for CVE-level discrimination
      (same-CVE = positive), which does not explicitly encourage
      CWE clustering. CWE silhouette remains negative for both
      GIN variants.
    \item \textbf{Combined still has higher absolute quality.}
      The multi-view Combined embedder benefits from complementary
      spectral (NetLSD) and colour-refinement (WL) signals,
      achieving the best silhouette and kNN purity overall.
      However, its intra/inter ratio (0.685) is \emph{worse} than
      trained GIN-Struct (0.305), indicating that the trained model
      has tighter CVE clusters despite lower absolute separability.
  \end{itemize}

  \subsubsection{Fusion with Trained GIN-Struct}
  \label{sec:gin-struct-fusion}

  The central question: \emph{does replacing the frozen random GIN
  with the trained GIN-Struct in the \texttt{norm\_concat\_pca}
  fusion strategy improve retrieval?}

  We compare the two fusion variants head-to-head: both use
  identical NetLSD and WL sub-embedders, identical L2-normalisation
  and PCA pipeline, and differ only in the GIN component.

  \begin{table}[ht]
  \centering
  \caption{Fusion comparison: frozen GIN vs.\ trained GIN-Struct
    in the \texttt{norm\_concat\_pca} pipeline. 225~index,
    73~queries, seed\,=\,42.}
  \label{tab:gin-struct-fusion}
  \small
  \renewcommand{\arraystretch}{1.2}
  \begin{tabular}{l cccc}
    \toprule
    \textbf{GIN component in fusion}
      & \textbf{hit@1} & \textbf{hit@5} & \textbf{hit@10}
      & \textbf{MRR} \\
    \midrule
    Frozen GIN (random)
      & 86.3\% & 97.3\% & 97.3\% & 0.907 \\
    \textbf{Trained GIN-Struct}
      & \textbf{89.0\%} & \textbf{98.6\%} & \textbf{98.6\%}
      & \textbf{0.929} \\
    \midrule
    \textit{Improvement}
      & \textit{+2.7pp} & \textit{+1.4pp} & \textit{+1.4pp}
      & \textit{+0.022} \\
    \bottomrule
  \end{tabular}
  \end{table}

  \subsubsection{Analysis}

  \paragraph{Training compounds with fusion.}
  The trained GIN-Struct improves the fusion from 86.3\% to 89.0\%
  hit@1 ($+2.7$~pp) and MRR from 0.907 to 0.929.
  This demonstrates that the discriminative structure learned via
  triplet loss is \emph{complementary} to NetLSD and WL---it
  captures CVE-specific topology patterns that the spectral and
  colour-refinement descriptors miss.

  \paragraph{Diminishing returns at high performance.}
  The absolute improvement ($+2.7$~pp) is smaller than the
  standalone improvement ($+5.5$~pp) because the fusion already
  achieves high accuracy via complementary sub-embedders.
  At 89.0\% hit@1, only 8 out of 73 queries fail---these represent
  genuinely difficult cases (e.g.\ NULL-pointer-dereference
  patterns that are structurally near-identical across unrelated
  CVEs).

  \paragraph{Near-perfect hit@5.}
  The 98.6\% hit@5 means only 1 query out of 73 fails to have the
  correct CVE in its top-5 neighbours.
  This is relevant for a production RAG system where
  re-ranking the top-5 candidates (e.g.\ via CodeBERT similarity
  or an LLM judge) is feasible.

  \paragraph{New best result.}
  The \texttt{norm\_concat\_pca} with trained GIN-Struct achieves
  \textbf{89.0\% hit@1 / 0.929 MRR}---the strongest retrieval
  result in the project to date.
  This is achieved entirely from structural features (node types
  + graph topology), without any code text or LLM-based embeddings.

  % ─────────────────────────────────────────────────────────────────
  \subsection{Learnings and Takeaways}
  \label{sec:learnings}

  \begin{enumerate}[leftmargin=*, label=\textbf{\arabic*.}]

    \item \textbf{Normalisation before fusion is critical.}
      The single most impactful change across all experiments was
      L2-normalising sub-embedders before concatenation.
      Without it, PCA is dominated by whichever sub-embedder has
      the largest scale, effectively discarding complementary signals.
      This yielded a $+16$~pp improvement with zero computational cost.

    \item \textbf{Structural features suffice for vulnerability retrieval.}
      The best result (89.0\% hit@1) uses only node types and graph
      topology---no code text, no LLM, no pre-trained language model.
      This has practical implications: the system is language-agnostic
      (works on any language with a CPG exporter), does not require
      GPU inference for CodeBERT, and produces deterministic embeddings.

    \item \textbf{Warm-starting from frozen weights is effective.}
      Initialising the trainable GIN from the frozen random GIN
      (exact architecture match, cosine sim 1.0 at init) allows
      triplet fine-tuning to converge quickly (94 epochs, 8~min)
      while preserving the Johnson--Lindenstrauss distance properties
      of the random initialisation as a regulariser.

    \item \textbf{BatchNorm mode matters for GIN.}
      When using BatchNorm in a GIN, the frozen model's embeddings
      use per-batch statistics (since running stats are never updated).
      At inference time, the trained model must also use
      \texttt{model.train()} mode to match per-batch statistics,
      rather than the standard \texttt{model.eval()} which uses
      (uninitialised) running statistics.
      This is a subtle but critical implementation detail.

    \item \textbf{CVE-level triplet loss does not transfer to CWE clustering.}
      The trained model dramatically improves CVE-level metrics
      (intra/inter ratio drops from 0.538 to 0.305) but CWE-level
      silhouette remains near-random ($-0.31$).
      CWE-level retrieval would require CWE-aware labels or a
      hierarchical loss.

    \item \textbf{Complementary views compound.}
      Each sub-embedder captures different structural aspects:
      NetLSD captures global spectral shape, WL captures local
      neighbourhood patterns via colour refinement, and GIN captures
      multi-hop message-passing topology.
      Their fusion consistently outperforms any individual component,
      and improving one component (GIN-Struct) lifts the fusion
      without degrading others.

    \item \textbf{The agent gap is real and quantifiable.}
      A 17.8~pp retrieval gap translates almost linearly into a
      15~pp patch quality (CodeBLEU proxy) loss.
      Improving retrieval from 68.5\% to 89.0\% hit@1 (+20.5~pp)
      should yield approximately $+18$~pp in patch quality---a
      substantial improvement that justifies continued investment
      in retrieval quality.

    \item \textbf{NULL-pointer-dereference is the hardest class.}
      Across all experiments, NPD vulnerabilities consistently
      underperform because their fixes (adding a null check) are
      structurally near-identical across unrelated CVEs.
      Disambiguating these requires function-level context
      (e.g.\ function name, call site) that current topology-only
      embedders do not capture.

  \end{enumerate}

  % ─────────────────────────────────────────────────────────────────
  \subsection{Summary of Results}
  \label{sec:results-summary}

  Table~\ref{tab:results-summary} consolidates key results
  across all experiments.

  \begin{table}[ht]
  \centering
  \caption{Consolidated retrieval results across experiments.
    All on the same 225/73 index/query split, $G_{\mathrm{vuln}}$,
    seed\,=\,42.}
  \label{tab:results-summary}
  \small
  \renewcommand{\arraystretch}{1.2}
  \begin{tabular}{l l cccc}
    \toprule
    \textbf{Experiment} & \textbf{Configuration}
      & \textbf{hit@1} & \textbf{hit@5}
      & \textbf{MRR} \\
    \midrule
    Individual & Frozen GIN
      & 72.6\% & 87.7\% & 0.788 \\
    Individual & Trained GIN-Struct
      & 78.1\% & 93.2\% & 0.826 \\
    \midrule
    Fusion & concat\_pca (raw)
      & 64.4\% & 79.5\% & 0.707 \\
    Fusion & norm\_concat\_pca (frozen)
      & 86.3\% & 97.3\% & 0.907 \\
    Fusion & norm\_concat\_pca (trained GIN-Struct)
      & \textbf{89.0\%} & \textbf{98.6\%} & \textbf{0.929} \\
    \midrule
    Agent & Embedding retriever (patch quality)
      & 68.5\% & 80.8\% & 0.733 \\
    Agent & Oracle retriever (patch quality)
      & 86.3\% & 90.4\% & 0.880 \\
    \midrule
    Verification & Diff-based patch verification
      & 73.6\% & 90.3\% & --- \\
    \bottomrule
  \end{tabular}
  \end{table}

  % ─────────────────────────────────────────────────────────────────
  \subsection{Next Steps}
  \label{sec:next-steps}

  Based on these findings, the following directions are prioritised:

  \begin{enumerate}[leftmargin=*, label=\textbf{\arabic*.}]
    \item \textbf{Deploy trained GIN-Struct in the agent pipeline.}
      Replace the frozen GIN in the production retriever with the
      trained GIN-Struct checkpoint to realise the 89.0\% hit@1 in
      end-to-end patch generation.
    \item \textbf{CWE-aware training.}
      Add a hierarchical triplet loss that uses CWE class as a
      coarse label, encouraging same-CWE graphs to cluster even when
      they are different CVEs.
    \item \textbf{Function-name augmentation for NPD disambiguation.}
      Add function-name hashing or a lightweight code-token feature
      to break ties between structurally identical null-check patches.
    \item \textbf{Re-ranking with CodeBERT.}
      Use the top-5 candidates from the structural retriever and
      re-rank with CodeBERT code similarity to combine structural
      recall (98.6\% hit@5) with semantic precision.
    \item \textbf{Scale to full dataset.}
      Extend evaluation to the full BigVul ($\sim$3,500 CVEs) and
      CVEFixes ($\sim$10,000 pairs) to validate generalisation.
  \end{enumerate}

  % ─────────────────────────────────────────────────────────────────
  \section{CVEFixes: Cross-Vulnerability Family Clustering}
  \label{sec:cvefixes}

  All preceding experiments used the \textbf{AutoPatch} dataset
  (75~CVEs, same-CVE retrieval).
  This section evaluates a fundamentally different question on the
  \textbf{CVEFixes} dataset ($\sim$10\,000 method-level vulnerability
  fixes): \emph{can unsupervised embeddings of vulnerability code
  discriminate between vulnerability families (CWE classes) without
  labelled training?}

  If vulnerability families have distinct structural or textual
  signatures, an unsupervised embedding should produce clusters
  aligned with the ground-truth family labels---enabling
  CWE-aware retrieval without supervised fine-tuning.

  \subsection{Experimental Setup}

  \subsubsection{Dataset and Families}

  From CVEFixes, we select functions belonging to 5~vulnerability
  families (30~samples each, 150~candidates with fallbacks for
  Joern parse failures, yielding $\sim$94 valid graphs):

  \begin{center}
  \small
  \renewcommand{\arraystretch}{1.2}
  \begin{tabular}{ll}
    \toprule
    \textbf{Family} & \textbf{Representative CWEs} \\
    \midrule
    memory\_corruption & CWE-119, CWE-120, CWE-787 \\
    use\_after\_free   & CWE-416 \\
    integer\_issues    & CWE-190, CWE-191 \\
    null\_ptr\_deref   & CWE-476 \\
    race\_condition    & CWE-362, CWE-667 \\
    \bottomrule
  \end{tabular}
  \end{center}

  Each function is exported via Joern, producing a full CPG
  ($G_{\mathrm{before}}$) which is then sliced according to
  different strategies.

  \subsubsection{Slicing Strategies}

  Three graph slicing strategies are compared, each producing a
  subgraph that is then embedded and clustered:

  \begin{enumerate}[nosep]
    \item \textbf{Token-guided:} entry nodes are those whose
      \texttt{CODE} contains discriminative tokens identified by
      TF-IDF/SVM analysis of each vulnerability family (e.g.\
      \texttt{free}, \texttt{kfree} for use\_after\_free;
      \texttt{lock}, \texttt{mutex} for race\_condition;
      \texttt{NULL}, \texttt{null} for null\_ptr\_deref).
      2-hop BFS expansion along flow edges from entry nodes.
    \item \textbf{Diff-based:} entry nodes are changed nodes
      (diff\_weight $> 0.3$) identified by semantic graph diff
      between pre-patch and post-patch CPGs. 2-hop expansion.
    \item \textbf{Random:} entry nodes are randomly selected
      (size-matched to token-guided). Controls for whether any
      subgraph of similar size would produce similar clustering.
  \end{enumerate}

  \subsubsection{Embedders}

  Seven embedders are evaluated:

  \begin{itemize}[nosep]
    \item \textbf{GIN} --- frozen random-weight 3-layer GIN
      (structural baseline)
    \item \textbf{GIN-enriched} --- GIN with 17-dimensional input
      features (11 node types + 5 diff labels + 1 diff\_weight)
    \item \textbf{Combined} --- NetLSD + WL + GIN, PCA to 128-d
    \item \textbf{Combined-enriched} --- Combined with enriched
      GIN component
    \item \textbf{CodeBERT-pattern} --- 34-d structural vulnerability
      patterns + 768-d CodeBERT on changed code, PCA to 128-d
    \item \textbf{CodeBERT-flow} --- CodeBERT with graph-flow-ordered
      code sequencing (BFS along REACHING\_DEF/CFG edges from seed
      nodes, inspired by GraphFVD)
    \item \textbf{CodeBERT-flow-pattern} --- 34-d patterns + flow-ordered
      CodeBERT fusion
  \end{itemize}

  \subsubsection{Evaluation Protocol}

  For each (strategy, embedder) pair:
  \begin{enumerate}[nosep]
    \item Embed all valid graph slices ($n \approx 94$).
    \item K-means cluster with $k = 5$ (number of families).
    \item Compute clustering metrics vs.\ ground-truth family labels:
      ARI, Homogeneity, NMI.
    \item Compute retrieval metrics: for each sample as query,
      rank all others by cosine similarity and measure P@$k$, MAP,
      MRR.
  \end{enumerate}

  \subsection{Iteration 1: Baseline Structural Embedders}

  The initial experiment tested whether GIN and Combined embedders
  on token-guided slices produce family-aligned clusters.

  \begin{table}[ht]
  \centering
  \caption{Clustering results --- structural embedders.}
  \label{tab:cvefixes-structural}
  \small
  \renewcommand{\arraystretch}{1.2}
  \begin{tabular}{ll rrr}
    \toprule
    \textbf{Strategy} & \textbf{Embedder}
      & \textbf{ARI} & \textbf{Homog.} & \textbf{NMI} \\
    \midrule
    token\_guided & GIN             &  0.015 & 0.064 & 0.071 \\
    token\_guided & Combined        & $-$0.001 & 0.057 & 0.057 \\
    diff\_based   & GIN             &  0.029 & 0.108 & 0.111 \\
    diff\_based   & Combined        &  0.026 & 0.091 & 0.096 \\
    random        & GIN             & $-$0.011 & 0.073 & 0.080 \\
    random        & Combined        & $-$0.003 & 0.087 & 0.091 \\
    \bottomrule
  \end{tabular}
  \end{table}

  \textbf{Finding:} All structural embedders produce near-zero ARI.
  Graph topology alone does not discriminate between vulnerability
  families. The frozen random GIN and Combined embedders capture
  graph shape, but different vulnerability types do not have
  sufficiently distinct topological signatures to cluster.

  \subsection{Iteration 2: Enriched Structural Features}

  Hypothesising that diff-awareness might help, we created enriched
  variants (GIN-enriched, Combined-enriched) that encode
  diff status and weight as node features (17-dim input: 11~node
  types + 5~diff labels + 1~diff\_weight).

  \begin{table}[ht]
  \centering
  \caption{Enriched vs.\ baseline structural embedders.}
  \label{tab:cvefixes-enriched}
  \small
  \renewcommand{\arraystretch}{1.2}
  \begin{tabular}{ll rrr}
    \toprule
    \textbf{Strategy} & \textbf{Embedder}
      & \textbf{ARI} & \textbf{Homog.} & \textbf{NMI} \\
    \midrule
    token\_guided & GIN             &  0.015 & 0.064 & 0.071 \\
    token\_guided & GIN-enriched    & $-$0.013 & 0.037 & 0.046 \\
    token\_guided & Combined        & $-$0.001 & 0.057 & 0.057 \\
    token\_guided & Combined-enriched & $-$0.001 & 0.057 & 0.057 \\
    \bottomrule
  \end{tabular}
  \end{table}

  \textbf{Finding:} Enriched variants perform \emph{worse} than
  baselines. Untrained random-weight GNNs cannot leverage additional
  features---they only add noise dimensions to the projection.
  The extra features would require supervised training (as in
  GIN-Struct, Section~\ref{sec:gin-struct}) to be useful.

  \subsection{Iteration 3: CodeBERT-Pattern (Semantic + Structural Fusion)}

  Since topology is insufficient, we tested CodeBERT-pattern which
  fuses 34 hand-crafted vulnerability pattern features (data-flow
  reachability, boundary edges, diff composition) with 768-d
  CodeBERT representations of the changed code.

  \begin{table}[ht]
  \centering
  \caption{CodeBERT-pattern clustering results.}
  \label{tab:cvefixes-codebert}
  \small
  \renewcommand{\arraystretch}{1.2}
  \begin{tabular}{ll rrr}
    \toprule
    \textbf{Strategy} & \textbf{Embedder}
      & \textbf{ARI} & \textbf{Homog.} & \textbf{NMI} \\
    \midrule
    token\_guided & codebert\_pattern & 0.068 & 0.154 & 0.154 \\
    diff\_based   & codebert\_pattern & 0.017 & 0.099 & 0.100 \\
    random        & codebert\_pattern & 0.019 & 0.113 & 0.125 \\
    \bottomrule
  \end{tabular}
  \end{table}

  \textbf{Finding:} CodeBERT-pattern on token-guided slices achieves
  the best clustering (ARI\,=\,0.068, NMI\,=\,0.154), confirming
  that discriminative textual tokens combined with semantic
  embedding provides marginally more signal than topology alone.
  However, the improvement over random baseline (ARI\,$\approx$\,0)
  is negligible in absolute terms.

  \subsection{Iteration 4: Graph-Flow-Ordered Code Sequencing}

  Inspired by GraphFVD~\cite{graphfvd2024}, we hypothesised that
  reordering code fed to CodeBERT by graph traversal (following
  REACHING\_DEF $\to$ CFG $\to$ CDG edges from seed nodes via BFS)
  would allow the self-attention mechanism to implicitly capture
  data-flow relationships between adjacent tokens.

  \begin{table}[ht]
  \centering
  \caption{Flow-ordered vs.\ line-ordered CodeBERT sequencing.}
  \label{tab:cvefixes-flow}
  \small
  \renewcommand{\arraystretch}{1.2}
  \begin{tabular}{ll rrr}
    \toprule
    \textbf{Strategy} & \textbf{Embedder}
      & \textbf{ARI} & \textbf{Homog.} & \textbf{NMI} \\
    \midrule
    token\_guided & codebert\_pattern       & 0.068 & 0.154 & 0.154 \\
    token\_guided & codebert\_flow          & 0.011 & 0.076 & 0.077 \\
    token\_guided & codebert\_flow\_pattern & 0.050 & 0.115 & 0.116 \\
    \bottomrule
  \end{tabular}
  \end{table}

  \textbf{Finding:} Flow-ordered sequencing \emph{degrades}
  performance. BFS traversal from seed nodes jumps across distant
  parts of a function, producing scrambled text that breaks
  CodeBERT's syntactic expectations (it was pre-trained on
  line-ordered source code). Line-number ordering preserves the
  syntactic locality that CodeBERT relies on for meaningful
  representations.

  \subsection{Retrieval Metrics}

  To provide a more fine-grained evaluation beyond clustering, we
  computed retrieval metrics for the best configuration
  (token\_guided + codebert\_pattern). For each of the 94~samples,
  all others are ranked by cosine similarity and we measure
  whether same-family items are retrieved first.

  \begin{table}[ht]
  \centering
  \caption{Retrieval metrics for codebert\_pattern on token-guided
    slices ($n=94$, $k=5$ families). Random baseline:
    P@1\,$\approx$\,0.20, MAP\,$\approx$\,0.20.}
  \label{tab:cvefixes-retrieval}
  \small
  \renewcommand{\arraystretch}{1.2}
  \begin{tabular}{l c}
    \toprule
    \textbf{Metric} & \textbf{Value} \\
    \midrule
    P@1  & 0.330 \\
    P@3  & 0.305 \\
    P@5  & 0.262 \\
    P@10 & 0.237 \\
    MAP  & 0.257 \\
    MRR  & 0.532 \\
    \bottomrule
  \end{tabular}
  \end{table}

  \begin{table}[ht]
  \centering
  \caption{Per-family retrieval breakdown.}
  \label{tab:cvefixes-retrieval-family}
  \small
  \renewcommand{\arraystretch}{1.2}
  \begin{tabular}{l ccccc}
    \toprule
    \textbf{Family}
      & \textbf{P@1} & \textbf{P@3} & \textbf{P@5}
      & \textbf{MAP} & \textbf{MRR} \\
    \midrule
    integer\_issues     & 0.333 & 0.259 & 0.211 & 0.245 & 0.496 \\
    memory\_corruption  & 0.353 & 0.275 & 0.200 & 0.235 & 0.544 \\
    null\_ptr\_deref    & 0.150 & 0.250 & 0.250 & 0.243 & 0.424 \\
    race\_condition     & 0.550 & 0.517 & 0.420 & 0.300 & 0.732 \\
    use\_after\_free    & 0.263 & 0.211 & 0.211 & 0.255 & 0.460 \\
    \bottomrule
  \end{tabular}
  \end{table}

  \textbf{Finding:} Retrieval performance is only marginally above
  the random baseline (P@1\,=\,0.330 vs.\ random
  $\approx$\,0.20; MAP\,=\,0.257 vs.\ $\approx$\,0.20). The
  $\sim$13~pp lift is not actionable for a retrieval system.
  Race condition shows the strongest signal
  (P@1\,=\,0.550, MRR\,=\,0.732), likely due to distinctive
  lock/unlock patterns in the code text. Null pointer
  dereference is the weakest (P@1\,=\,0.150), consistent with
  earlier findings that null-check patterns are structurally
  homogeneous across unrelated CVEs.

  \subsection{Consolidated Results}

  Table~\ref{tab:cvefixes-all} presents the complete
  strategy$\times$embedder results grid.

  \begin{table}[ht]
  \centering
  \caption{Full clustering results: all strategy$\times$embedder
    combinations on CVEFixes ($n \approx 94$ per strategy).}
  \label{tab:cvefixes-all}
  \small
  \renewcommand{\arraystretch}{1.1}
  \setlength{\tabcolsep}{3.5pt}
  \begin{tabular}{ll rrr r}
    \toprule
    \textbf{Strategy} & \textbf{Embedder}
      & \textbf{ARI} & \textbf{Homog.} & \textbf{NMI}
      & $n$ \\
    \midrule
    token\_guided & GIN                    &  0.015 & 0.064 & 0.071 & 94 \\
    token\_guided & GIN-enriched           & $-$0.013 & 0.037 & 0.046 & 94 \\
    token\_guided & Combined               & $-$0.001 & 0.057 & 0.057 & 94 \\
    token\_guided & Combined-enriched      & $-$0.001 & 0.057 & 0.057 & 94 \\
    token\_guided & codebert\_pattern      &  \textbf{0.068} & \textbf{0.154} & \textbf{0.154} & 94 \\
    token\_guided & codebert\_flow         &  0.011 & 0.076 & 0.077 & 94 \\
    token\_guided & codebert\_flow\_pattern &  0.050 & 0.115 & 0.116 & 94 \\
    \midrule
    diff\_based   & GIN                    &  0.029 & 0.108 & 0.111 & 92 \\
    diff\_based   & GIN-enriched           &  0.038 & 0.090 & 0.093 & 92 \\
    diff\_based   & Combined               &  0.026 & 0.091 & 0.096 & 92 \\
    diff\_based   & Combined-enriched      &  0.035 & 0.101 & 0.109 & 92 \\
    diff\_based   & codebert\_pattern      &  0.017 & 0.099 & 0.100 & 92 \\
    diff\_based   & codebert\_flow         &  0.018 & 0.094 & 0.094 & 92 \\
    diff\_based   & codebert\_flow\_pattern & $-$0.002 & 0.074 & 0.074 & 92 \\
    \midrule
    random        & GIN                    & $-$0.011 & 0.073 & 0.080 & 70 \\
    random        & GIN-enriched           & $-$0.018 & 0.075 & 0.080 & 70 \\
    random        & Combined               & $-$0.003 & 0.087 & 0.091 & 70 \\
    random        & Combined-enriched      & $-$0.003 & 0.087 & 0.091 & 70 \\
    random        & codebert\_pattern      &  0.019 & 0.113 & 0.125 & 70 \\
    random        & codebert\_flow         &  0.024 & 0.114 & 0.124 & 70 \\
    random        & codebert\_flow\_pattern &  0.019 & 0.107 & 0.118 & 70 \\
    \bottomrule
  \end{tabular}
  \end{table}

  \subsection{Conclusions}
  \label{sec:cvefixes-conclusions}

  \begin{enumerate}[leftmargin=*, label=\textbf{\arabic*.}]

    \item \textbf{Unsupervised CWE family clustering does not work.}
      The best configuration (token\_guided + codebert\_pattern)
      achieves ARI\,=\,0.068 and MAP\,=\,0.257---barely above
      random chance ($\approx$\,0.20). No embedder, slicing strategy,
      or code ordering produces meaningful family-level clusters
      without supervised training.

    \item \textbf{Vulnerability type is not the dominant signal.}
      The embedding space is dominated by function-specific and
      project-specific variation (code style, naming conventions,
      function complexity) that overwhelms the subtle structural
      differences between CWE families. Two null-pointer-dereference
      fixes in different subsystems are more different from each other
      than from a use-after-free fix in the same subsystem.

    \item \textbf{Graph topology alone is uninformative for CWE
      discrimination.}
      All GIN/Combined variants (including enriched) produce
      ARI\,$\leq$\,0.038. The topological ``shape'' of a
      vulnerability slice does not encode which \emph{type} of
      vulnerability it is. This contrasts sharply with the AutoPatch
      results where structural embedders achieve 89\% hit@1---but
      that task is same-CVE retrieval (matching identical structure),
      not cross-CVE family discrimination.

    \item \textbf{Token-guided slicing helps only with CodeBERT.}
      The discriminative tokens (free, lock, null, etc.) provide
      useful signal only when processed by a language model that
      understands their semantics. For structural embedders, the
      token-guided slice is no better than random.

    \item \textbf{Graph-flow ordering is counterproductive for CodeBERT.}
      Reordering code by BFS over flow edges destroys the syntactic
      locality that CodeBERT expects from pre-training. The model
      was trained on source code in line order; presenting scrambled
      statement sequences degrades representation quality.

    \item \textbf{Race conditions are the most retrievable family.}
      P@1\,=\,0.550 and MRR\,=\,0.732 for race\_condition---likely
      because lock/unlock API calls (\texttt{mutex\_lock},
      \texttt{spin\_lock}) create distinctive textual patterns
      recognisable by CodeBERT. This suggests that vocabulary-based
      retrieval (keyword matching) may outperform embedding-based
      retrieval for CWE discrimination.

    \item \textbf{Implications for the RAG system.}
      Cross-CWE retrieval requires either:
      (a)~supervised contrastive learning with CWE labels
      (as demonstrated effective for same-CVE retrieval in
      Section~\ref{sec:gin-struct}), or
      (b)~a fundamentally different approach such as
      specification-level matching (comparing vulnerability
      descriptions or fix-lists rather than raw code embeddings).
      The current unsupervised embedding pipeline is effective for
      same-CVE retrieval (89\% hit@1) but cannot generalise to
      family-level discrimination.

  \end{enumerate}

  % ─────────────────────────────────────────────────────────────────
  \begin{thebibliography}{9}

  \bibitem{beyer1999nearest}
    K.~Beyer, J.~Goldstein, R.~Ramakrishnan, and U.~Shaft,
    ``When is `nearest neighbor' meaningful?''
    in \emph{Proc.\ 7th International Conference on Database Theory
    (ICDT)}, pp.~217--235, Springer, 1999.

  \bibitem{jolliffe2016pca}
    I.\,T.~Jolliffe and J.~Cadima,
    ``Principal component analysis: a review and recent developments,''
    \emph{Philosophical Transactions of the Royal Society~A},
    vol.~374, no.~2065, p.~20150202, 2016.

  \bibitem{raunak2019effective}
    V.~Raunak, V.~Gupta, and F.~Metze,
    ``Effective dimensionality reduction for word embeddings,''
    in \emph{Proc.\ 4th Workshop on Representation Learning for NLP
    (RepL4NLP)}, pp.~235--243, ACL, 2019.

  \bibitem{autopatch2024}
    J.~Park, J.~Lee, and S.~Ryu,
    ``AutoPatch: Automated Generation of Hotpatches for Real-Time
    Embedded Devices,''
    \emph{arXiv preprint arXiv:2404.xxxxx}, 2024.

  \bibitem{papineni2002bleu}
    K.~Papineni, S.~Roukos, T.~Ward, and W.-J.~Zhu,
    ``BLEU: a method for automatic evaluation of machine translation,''
    in \emph{Proc.\ 40th Annual Meeting of the Association for
    Computational Linguistics (ACL)}, pp.~311--318, 2002.

  \bibitem{ren2020codebleu}
    S.~Ren, D.~Guo, S.~Lu, L.~Zhou, S.~Liu, D.~Tang, N.~Sundaresan,
    M.~Zhou, A.~Blanco, and S.~Ma,
    ``CodeBLEU: a method for automatic evaluation of code synthesis,''
    \emph{arXiv preprint arXiv:2009.10297}, 2020.

  \bibitem{xia2023automated}
    C.\,S.~Xia, Y.~Wei, and L.~Zhang,
    ``Automated program repair in the era of large pre-trained
    language models,''
    in \emph{Proc.\ 45th International Conference on Software
    Engineering (ICSE)}, pp.~1482--1494, IEEE, 2023.

  \bibitem{graphfvd2024}
    Y.~Li, S.~Wang, and T.\,N.~Nguyen,
    ``GraphFVD: Fine-Grained Vulnerability Detection via Graph-Based
    Code Representation,''
    \emph{IEEE Transactions on Software Engineering}, 2024.

  \end{thebibliography}

  \end{document}